\clearpage\SourceLeaf{402}\ArithmeticHead{395}{R}{BEMERKUNGEN ZUR REIHENLEHRE. ABSCHNITT II.}
\begin{center}BEMERKUNGEN.\end{center}
\noindent\hspace*{\parindent}Die Begriffe und Sätze, die \textsc{Gauss} in dem vorstehenden Bruchstück entwickelt, sind lange vor der ersten Veröffentlichung dieser Aufzeichnung (1912 im Heft III der \textit{Materialien zu einer wissenschaftlichen Biographie von Gauss}, S. 136) von andern Mathematikern aufgestellt und bekannt gemacht worden. Was \textsc{Gauss} kleinste obere bezw. größte untere Grenze nennt, wird gewöhnlich auf B. \textsc{Bolzano} zurückgeführt\textsuperscript{*}); nach M. \textsc{Pasch}\textsuperscript{**}) sagt man dafür obere bezw. untere \textls[120]{Schranke}. Was \textsc{Gauss} als letzte obere bezw. letzte untere Grenze bezeichnet, findet sich bei \textsc{Cauchy} als la plus grande bezw. la plus petite des limites\textsuperscript{***}); P. \textsc{du Bois Reymond}\textsuperscript{\dag}) sagt dafür obere und untere \textls[120]{Unbestimmtheitsgrenze}. Der Begriff der obern Schranke findet sich übrigens schon im art. 6. Absatz 4. von \textsc{Gauss}' \textit{Demonstratio nova} (1799), Werke III, S. 10. Während letzte obere und untere Grenzen stets vorhanden sind, bildet das Vorhandensein einer \textls[120]{absoluten} Grenze (art. 5.), d. h. des Grenzwertes oder Limes für eine Reihe einen besondern Fall.

Die \textls[120]{Reihen}, die \textsc{Gauss} hier betrachtet und auf die er sich \textls[120]{absichtlich} beschränkt, würden mit einem neuern Kunstausdruck als abzählbare Mengen zu bezeichnen sein, die überdies in einer bestimmten Anordnung vorgelegt sind. Von einer Ausdehnung der Untersuchung auf den »Inbegriff einer jeden Anzahl beliebiger Größen«, d. h. also auf allgemeine lineare Punktmengen, versprach \textsc{Gauss} sich »wenig Nutzen« (siehe die einleitenden Worte, S. 390).

Zum art. 4., S. 393 ist noch zu bemerken, daß die Glieder der Reihe [1] die Koeffizienten sind in der Entwicklung der rationalen Funktion [2] nach positiven ganzen Potenzen von \(x\). Mit den aus der Entwicklung rationaler Funktionen entspringenden sogenannten \textls[120]{rekurrierenden Reihen}, besonders mit den zugehörigen \textls[120]{Relationsskalen} und ihren Verallgemeinerungen hat sich \textsc{Gauss} in jungen Jahren vielfach beschäftigt, siehe die Tagebuchaufzeichnungen Nr. 8 vom 26. Mai, Nr. 10 vom 3. Juni und Nr. 20 vom 16. Juli 1796.

Die letzte Seite der Handschrift enthält Rechnungen zur Osterformel und zwar mit Beispielen für das achtzehnte Jahrhundert; mit Rücksicht auf die Tagebuchnotiz Nr. 107 vom 16. Mai 1800 wird man daher die Abfassungszeit dieses Bruchstücks auf die Wende des XVIII. zum XIX. Jahrhundert ansetzen können.
\begin{flushright}\textsc{Schlesinger}.\end{flushright}
\par\medskip\noindent\rule{.24\linewidth}{.4pt}\par\smallskip
{\small
\noindent\hspace*{\parindent}\textsuperscript{*}) \textit{Rein analytischer Beweis des Lehrsatzes, dass zwischen je zwei Werthen, die ein entgegengesetztes Resultat gewähren, wenigstens eine reelle Wurzel der Gleichung liege}, Prag 1817, S. 41, vergl. O. \textsc{Stolz} Mathematische Annalen \textbf{18}, 1881, S. 257.
\par\noindent\hspace*{\parindent}\textsuperscript{**}) Mathematische Annalen \textbf{30}, 1887, S. 132.
\par\noindent\hspace*{\parindent}\textsuperscript{***}) \textit{Analyse algébrique}, 1821, S. 132, 151 bezw. 390, 399; Oeuvres, 2. série, t. III, S. 121, 136 bezw. 322, 329.
\par\noindent\hspace*{\parindent}\textsuperscript{\dag}) Antritts-Programm der Universität Freiburg, 1871, S. 3, vergl. \textit{Allgemeine Functionentheorie} I, 1882, S. 266.
\par}
\begin{center}\rule{.24\linewidth}{.4pt}\end{center}
\noindent\hfill50\textsuperscript{*}
